\chapter{Conclusão}\label{cap:conclusao}

A documentação relaciona o domínio de comparação de preços farmacêuticos aos requisitos, casos de uso, classes, esquema relacional, interfaces e planejamento de testes do ``Preço Bão''. A implementação consultada contém componentes de catálogo, autenticação, favoritos e observações históricas de preços, organizados em backend FastAPI, frontend React Native/Expo e persistência PostgreSQL. O DER e o dicionário tornam explícitas as tabelas associativas e os atributos herdados da autenticação.

A conferência também evidencia diferenças entre metas e implementação. Recuperação de senha no fluxo público, cálculo de distância na pesquisa, retenção automática de históricos por 12 meses e bloqueio de alterações por verificações de qualidade não foram demonstrados pelas fontes consultadas. As metas de desempenho, funcionamento em dispositivos móveis e conformidade de privacidade exigem validações específicas; a presença de código ou configuração não comprova esses resultados.

As Práticas 1 e 2 acrescentam uma análise de riscos de manutenção e um planejamento de intervenções. A primeira relaciona problemas de coleta, ambiente, banco e dependências a alternativas e estimativas. A segunda classifica quatro demandas, organiza prioridades e sprints e detalha coletivamente a sincronização de favoritos. A contribuição comprovada desses registros é documental: não há resultados de implantação, indicadores medidos ou execução de sprints a atribuir às propostas.

O plano de testes mantém os casos pendentes e suas datas originais. Assim, não se atribuem ao projeto resultados experimentais, métricas ou responsabilidades executadas sem evidência. A consolidação técnica fornece rastreabilidade para continuidade do desenvolvimento, mas a avaliação das entregas acadêmicas de cada período depende dos enunciados, produtos e diários que documentem sua realização.
